% Member 2 - final report, II.C Style Augmentation for Image Classification
\subsection{Style Augmentation for Image Classification}
\label{sec:style_aug_classification}

Using generated images as training data is not new. The ``domain
generalization through data augmentation'' line of work
\cite{zhou2023dgsurvey} has established that synthetic variation can close a
domain gap without touching the model, and several groups have
specialised the idea to appearance. The most direct precedent is style
augmentation itself~\cite{jackson2019style}, which randomises the style of
the training images so that the classifier cannot use it as a cue. Other
groups have perturbed the \emph{features} inside the network, without
generating any image at all. MixStyle~\\cite{zhou2021mixstyle} is the clearest
example of the second: it mixes the channel-wise mean and variance of a
feature map with the statistics of another sample, which removes the
domain-specific style of the input in feature space. This distinction
matters for our study, because the second strategy does not produce images
and therefore cannot be analysed as images. We include MixStyle as a
published baseline (Table~\ref{tab:main_results}) precisely because it is
the strongest competitor of the same family: it targets the same problem,
removes the same nuisance, and needs no generator.

The generator-based strategy is what this project studies, because it lets
the strength of the stylisation be an explicit, inspectable quantity. It
also lets us ask a question that feature-space methods cannot: \emph{how
much} style should be applied, and \emph{where} should it be applied? The
first question is about the total amount of appearance change, the second
about its spatial distribution. Our experiments separate them with two
controls that hold one constant while the other varies, which is the design
described in Section~\ref{sec:strength_relation}.

The prevailing intuition is that stronger stylisation is better, on the
grounds that it removes more of the domain-specific appearance cue. Our
results qualify that intuition. The accuracy of a styled training set is
governed by its average style strength, and once that is matched, the
spatial distribution of the strength contributes nothing measurable
(Section~\ref{sec:strength_relation}). Reporting the strengths explicitly
is therefore not a detail of the method but a precondition for a fair
comparison, and we report it for every configuration in
Table~\ref{tab:main_results}.
